
\section{Architectural Overview}\label{sec:irie-arch}

The framework is organized as a four‑layer interface stack that
flows from user‑facing web pages down to operating‑system resources. 
Each layer holds a distinct responsibility and
communicates with its neighbors solely through deliberately thin,
technology‑neutral interfaces. 
This partitioning both localizes change %---new database engines or numerical solvers can be introduced without disturbing the presentation tier---
and establishes clear trust boundaries for security and audit.

\subsection{Presentation Layer}\label{presentation-layer}

The topmost layer delivers HTML dashboards, REST endpoints, and
programmatic client libraries \citep{richardson2015restful}. 
Its single obligation is to translate internal domain concepts into artifacts that humans or external services can consume. 
All interaction with the analysis core is mediated through
type‑checked view models; raw database cursors or solver handles never cross this boundary. 
The separation prevents tightly coupled
dependencies and allows rapid evolution of the interface without risking regressions in numerical behavior.

\subsection{Orchestration Layer}\label{orchestration-layer}

Beneath the presentation tier sits a scheduler that governs the
life‑cycle of monitoring workflows. 
It listens for \emph{Event} notifications, allocates compute resources, dispatches \emph{Predictor}
instances, and reconciles their outputs into an \emph{Evaluation}.
Orchestration is intentionally stateless between jobs. 
Every mutable artifact is persisted by the Domain Layer so that workflows can be replayed verbatim
for audit or research replication. 
Because the scheduler interacts with
predictors only through a minimal launch protocol, additional numerical
engines can be registered without changes to scheduling code.

\subsection{Domain Layer}\label{domain-layer}

The Domain Layer captures the five abstractions that underpin the
monitoring logic. These objects expose behavior through stable
interfaces while hiding implementation details such as schema migration,
file‑system layout, or communication middleware. 
Two design decisions reinforce the layer's longevity:

\begin{itemize}
\tightlist
\item
  \textbf{Composition over specialization.} New predictor capabilities
  are introduced by assembling existing behavior blocks rather than by
  extending deep inheritance hierarchies. This approach limits
  cross‑cutting side‑effects and keeps object graphs shallow.
\item
  \textbf{Strict data immutability at workflow boundaries.} An
  \emph{Event} record is never modified after creation; an
  \emph{Evaluation} becomes read‑only once published. Immutability
  simplifies concurrent access, supports content‑addressable caching,
  and eliminates entire classes of race conditions.
\end{itemize}

\subsection{Infrastructure Layer}\label{infrastructure-layer}

The bottom layer abstracts operating‑system concerns: database
connectivity, object storage, job queues, and low‑level telemetry. These
services are accessed through narrow gateway classes so that the
remainder of the codebase remains indifferent to whether the platform
runs on a developer laptop or a cloud cluster. Changes in deployment
target therefore require only the replacement of gateway adapters, not a
refactor of business logic.

\subsection{Layer Collaboration}\label{layer-collaboration}

Communication proceeds strictly downward. For example, when the
Presentation Layer receives a request for the latest bridge health
report it calls the Orchestration Layer, which in turn queries the
Domain Layer for immutable \emph{Evaluation} objects. 
Only the Infrastructure Layer is permitted to reach outside the codebase. 
Upward calls are prohibited; infrastructure code does not emit direct UI updates, and domain objects have no
knowledge of HTTP\footnote{The Hypertext transfer protocol is a standard protocol established by \cite{fielding2022http}.} semantics. 
The unidirectional rule guards against hidden coupling and ensures that each layer can be tested in isolation.

\section{Prediction Subsystem}\label{sec:scheduling}

The design objective of the Prediction Subsystem is to permit the
continuous integration of new modeling techniques while shielding the
rest of the platform from solver‑specific details. 
To this end, each
concrete predictor is packaged behind a \textbf{Strategy} object \citep{gamma2011design} that
realizes a simplified four‑phase template: \emph{initialize},
\emph{execute}, \emph{harvest}, and \emph{summarize}. 
A lightweight
\textbf{Factory Registry} maps protocol identifiers to concrete strategy
instances at runtime. 
Adding a novel predictor therefore requires no modification to orchestration code, only the registration of a new
strategy class and, if desired, the definition of new metric types to
reflect its particular output.

During execution a predictor writes its raw results to a unique run
directory and emits a structured summary describing the provenance of
every number that ultimately appears on a dashboard. This provenance
trail is critical to after‑action audits and mirrors the reproducibility
guarantees advocated in earlier software for finite‑element research.
The same provenance mechanism underpins a run‑level caching scheme: if
an identical event--predictor combination is encountered twice, the
second request can be satisfied by replaying the cached results instead
of performing a new analysis.

\subsection{Metrics}\label{metrics}

Health metrics are immutable value objects \citep{fowler2013patterns} that bind three pieces of information: a \emph{label} that identifies the physical
quantity (e.g., fundamental period, peak deck acceleration, column
strain state), a compact payload containing the numerical
value(s), and a \emph{verdict} that classifies the value against user‑defined thresholds.
Metrics understand how to render themselves in the web interface, how to serialize into archival formats, and how to participate in evaluation‑level aggregation. 
Because all predictors emit metrics rather than solver‑specific artifacts, downstream consumers need not know whether a given metric was extracted from a finite‑element response history,
inferred by a statistical emulator, or read from a fragility table.

\section{Event and Evaluation Workflow}\label{sec:policy}

\begin{figure}[ht]
    \centering
    \resizebox{0.9\linewidth}{!}{% generated by Plantuml 1.2025.4       
\definecolor{plantucolor0000}{RGB}{0,0,0}
\definecolor{plantucolor0001}{RGB}{255,255,255}
\definecolor{plantucolor0002}{RGB}{24,24,24}
\definecolor{plantucolor0003}{RGB}{227,227,227}
\definecolor{plantucolor0004}{RGB}{238,238,238}
\begin{tikzpicture}[yscale=-1
,pstyle0/.style={color=black,line width=1.5pt}
,pstyle1/.style={color=plantucolor0001,line width=0.0pt}
,pstyle2/.style={color=plantucolor0002,line width=0.5pt,dash pattern=on 5.0pt off 5.0pt}
,pstyle3/.style={color=plantucolor0002,fill=plantucolor0003,line width=0.5pt}
,pstyle4/.style={color=plantucolor0002,fill=plantucolor0002,line width=1.0pt}
,pstyle5/.style={color=plantucolor0002,line width=1.0pt}
]
\draw[pstyle0] (112.14pt,140pt) rectangle (400.44pt,247pt);
\draw[pstyle1] (38.45pt,40pt) rectangle (46.45pt,312pt);
\draw[pstyle2] (42pt,40pt) -- (42pt,312pt);
\draw[pstyle1] (172.23pt,40pt) rectangle (180.23pt,312pt);
\draw[pstyle2] (176.14pt,40pt) -- (176.14pt,312pt);
\draw[pstyle1] (278.19pt,40pt) rectangle (286.19pt,312pt);
\draw[pstyle2] (281.975pt,40pt) -- (281.975pt,312pt);
\draw[pstyle1] (391.44pt,40pt) rectangle (399.44pt,312pt);
\draw[pstyle2] (395.285pt,40pt) -- (395.285pt,312pt);
\draw[pstyle1] (439.955pt,40pt) rectangle (447.955pt,312pt);
\draw[pstyle2] (443.595pt,40pt) -- (443.595pt,312pt);
\draw[pstyle3] (5pt,20pt) arc (180:270:5pt) -- (10pt,15pt) -- (74.9pt,15pt) arc (270:360:5pt) -- (79.9pt,20pt) -- (79.9pt,34pt) arc (0:90:5pt) -- (74.9pt,39pt) -- (10pt,39pt) arc (90:180:5pt) -- (5pt,34pt) -- cycle;
\node at (12pt,22pt)[below right,color=black,inner sep=0]{Web Gateway};
\draw[pstyle3] (5pt,316pt) arc (180:270:5pt) -- (10pt,311pt) -- (74.9pt,311pt) arc (270:360:5pt) -- (79.9pt,316pt) -- (79.9pt,330pt) arc (0:90:5pt) -- (74.9pt,335pt) -- (10pt,335pt) arc (90:180:5pt) -- (5pt,330pt) -- cycle;
\node at (12pt,318pt)[below right,color=black,inner sep=0]{Web Gateway};
\draw[pstyle3] (122.14pt,10pt) arc (180:270:5pt) -- (127.14pt,5pt) -- (225.32pt,5pt) arc (270:360:5pt) -- (230.32pt,10pt) -- (230.32pt,34pt) arc (0:90:5pt) -- (225.32pt,39pt) -- (127.14pt,39pt) arc (90:180:5pt) -- (122.14pt,34pt) -- cycle;
\node at (146.31pt,12pt)[below right,color=black,inner sep=0]{Orchestration};
\node at (129.14pt,22pt)[below right,color=black,inner sep=0]{(Evaluation Daemon)};
\draw[pstyle3] (122.14pt,316pt) arc (180:270:5pt) -- (127.14pt,311pt) -- (225.32pt,311pt) arc (270:360:5pt) -- (230.32pt,316pt) -- (230.32pt,340pt) arc (0:90:5pt) -- (225.32pt,345pt) -- (127.14pt,345pt) arc (90:180:5pt) -- (122.14pt,340pt) -- cycle;
\node at (146.31pt,318pt)[below right,color=black,inner sep=0]{Orchestration};
\node at (129.14pt,328pt)[below right,color=black,inner sep=0]{(Evaluation Daemon)};
\draw[pstyle3] (248.975pt,20pt) arc (180:270:5pt) -- (253.975pt,15pt) -- (310.405pt,15pt) arc (270:360:5pt) -- (315.405pt,20pt) -- (315.405pt,34pt) arc (0:90:5pt) -- (310.405pt,39pt) -- (253.975pt,39pt) arc (90:180:5pt) -- (248.975pt,34pt) -- cycle;
\node at (255.975pt,22pt)[below right,color=black,inner sep=0]{Predictor [i]};
\draw[pstyle3] (248.975pt,316pt) arc (180:270:5pt) -- (253.975pt,311pt) -- (310.405pt,311pt) arc (270:360:5pt) -- (315.405pt,316pt) -- (315.405pt,330pt) arc (0:90:5pt) -- (310.405pt,335pt) -- (253.975pt,335pt) arc (90:180:5pt) -- (248.975pt,330pt) -- cycle;
\node at (255.975pt,318pt)[below right,color=black,inner sep=0]{Predictor [i]};
\draw[pstyle3] (371.285pt,10pt) arc (180:270:5pt) -- (376.285pt,5pt) -- (414.595pt,5pt) arc (270:360:5pt) -- (419.595pt,10pt) -- (419.595pt,34pt) arc (0:90:5pt) -- (414.595pt,39pt) -- (376.285pt,39pt) arc (90:180:5pt) -- (371.285pt,34pt) -- cycle;
\node at (378.285pt,12pt)[below right,color=black,inner sep=0]{Domain};
\node at (383.275pt,22pt)[below right,color=black,inner sep=0]{Layer};
\draw[pstyle3] (371.285pt,316pt) arc (180:270:5pt) -- (376.285pt,311pt) -- (414.595pt,311pt) arc (270:360:5pt) -- (419.595pt,316pt) -- (419.595pt,340pt) arc (0:90:5pt) -- (414.595pt,345pt) -- (376.285pt,345pt) arc (90:180:5pt) -- (371.285pt,340pt) -- cycle;
\node at (378.285pt,318pt)[below right,color=black,inner sep=0]{Domain};
\node at (383.275pt,328pt)[below right,color=black,inner sep=0]{Layer};
\draw[pstyle3] (429.595pt,20pt) arc (180:270:5pt) -- (434.595pt,15pt) -- (453.315pt,15pt) arc (270:360:5pt) -- (458.315pt,20pt) -- (458.315pt,34pt) arc (0:90:5pt) -- (453.315pt,39pt) -- (434.595pt,39pt) arc (90:180:5pt) -- (429.595pt,34pt) -- cycle;
\node at (436.595pt,22pt)[below right,color=black,inner sep=0]{DB};
\draw[pstyle3] (429.595pt,316pt) arc (180:270:5pt) -- (434.595pt,311pt) -- (453.315pt,311pt) arc (270:360:5pt) -- (458.315pt,316pt) -- (458.315pt,330pt) arc (0:90:5pt) -- (453.315pt,335pt) -- (434.595pt,335pt) arc (90:180:5pt) -- (429.595pt,330pt) -- cycle;
\node at (436.595pt,318pt)[below right,color=black,inner sep=0]{DB};
\draw[pstyle4] (164.23pt,73pt) -- (174.23pt,77pt) -- (164.23pt,81pt) -- (168.23pt,77pt) -- cycle;
\draw[pstyle5] (42.45pt,77pt) -- (170.23pt,77pt);
\node at (49.45pt,54pt)[below right,color=black,inner sep=0]{POST /events};
\node at (49.45pt,65pt)[below right,color=black,inner sep=0]{(sensor payload)};
\draw[pstyle4] (431.955pt,97pt) -- (441.955pt,101pt) -- (431.955pt,105pt) -- (435.955pt,101pt) -- cycle;
\draw[pstyle5] (176.23pt,101pt) -- (437.955pt,101pt);
\node at (183.23pt,89pt)[below right,color=black,inner sep=0]{insert Event};
\draw[pstyle4] (383.44pt,121pt) -- (393.44pt,125pt) -- (383.44pt,129pt) -- (387.44pt,125pt) -- cycle;
\draw[pstyle5] (176.23pt,125pt) -- (389.44pt,125pt);
\node at (183.23pt,113pt)[below right,color=black,inner sep=0]{query active Predictors};
\draw[color=black,fill=plantucolor0004,line width=1.5pt] (112.14pt,140pt) -- (178.54pt,140pt) -- (178.54pt,142pt) -- (168.54pt,152pt) -- (112.14pt,152pt) -- (112.14pt,140pt);
\draw[pstyle0] (112.14pt,140pt) rectangle (400.44pt,247pt);
\node at (127.14pt,141pt)[below right,color=black,inner sep=0]{\textbf{loop}};
\node at (193.54pt,142pt)[below right,color=black,inner sep=0]{\textbf{[for each Predictor]}};
\draw[pstyle4] (270.19pt,164pt) -- (280.19pt,168pt) -- (270.19pt,172pt) -- (274.19pt,168pt) -- cycle;
\draw[pstyle5] (176.23pt,168pt) -- (276.19pt,168pt);
\node at (183.23pt,156pt)[below right,color=black,inner sep=0]{launch(i, EventID)};
\draw[pstyle5] (282.19pt,192pt) -- (324.19pt,192pt);
\draw[pstyle5] (324.19pt,192pt) -- (324.19pt,205pt);
\draw[pstyle5] (283.19pt,205pt) -- (324.19pt,205pt);
\draw[pstyle4] (293.19pt,201pt) -- (283.19pt,205pt) -- (293.19pt,209pt) -- (289.19pt,205pt) -- cycle;
\node at (289.19pt,180pt)[below right,color=black,inner sep=0]{init→execute→harvest};
\draw[pstyle4] (187.23pt,235pt) -- (177.23pt,239pt) -- (187.23pt,243pt) -- (183.23pt,239pt) -- cycle;
\draw[color=plantucolor0002,line width=1.0pt,dash pattern=on 2.0pt off 2.0pt] (181.23pt,239pt) -- (281.19pt,239pt);
\node at (193.23pt,217pt)[below right,color=black,inner sep=0]{Metric set};
\node at (193.23pt,227pt)[below right,color=black,inner sep=0]{+ provenance};
\draw[pstyle4] (431.955pt,266pt) -- (441.955pt,270pt) -- (431.955pt,274pt) -- (435.955pt,270pt) -- cycle;
\draw[pstyle5] (176.23pt,270pt) -- (437.955pt,270pt);
\node at (183.23pt,258pt)[below right,color=black,inner sep=0]{insert Evaluation + Metrics};
\draw[pstyle4] (53.45pt,290pt) -- (43.45pt,294pt) -- (53.45pt,298pt) -- (49.45pt,294pt) -- cycle;
\draw[pstyle5] (47.45pt,294pt) -- (175.23pt,294pt);
\node at (59.45pt,282pt)[below right,color=black,inner sep=0]{notify(evaluation\_ready)};
\end{tikzpicture}
}
    \caption{Sequence diagram illustrating the process of evaluating an event.}
    \label{fig:irie-sequences}
\end{figure}

An event enters the system through a secure web service that
authenticates the request, stores raw sensor streams in object storage,
and inserts a new row in the \emph{Event} ledger. 
At that instant a long‑lived scheduler running in the
Orchestration Layer called the \textbf{Evaluation Daemon} receives a notification. 
The daemon queries the Asset record to discover which predictors are active and queues them for
execution according to resource‑aware policies. 
Predictors may run concurrently, and each operates in isolation within its own working directory.

When a predictor finishes, the daemon collects its metrics and inserts
them into the database under a new \emph{Evaluation} record. 
Once all predictors for an event have reported, the daemon invokes an
\textbf{Evaluation Policy} that scans the assembled metric set, assigns
an overall health state, and dispatches notifications to stakeholders.
The policy mechanism is again a Strategy: facilities managers can supply
their own decision logic, ranging from simple threshold checks to
Bayesian belief networks, without disturbing core services.

\Cref{fig:irie-sequences} illustrates the complete event path, from acceleration
time histories arriving at the web gateway to color‑coded health tiles
appearing on the bridge dashboard. 
Every arrow in the sequence diagram is traceable to a method on one of the abstractions.


\newcommand{\BlankLine}{}

\begin{algorithm}[h]
\caption{\enskip Top–level evaluation routine.}\label{alg:evaluate}
\begin{algorithmic}
  \Function{Evaluate}{$\mathit{event},\,\mathit{evaluation}$}
    \State $c_{\mathrm{ok}} \gets 0$
    \State $\mathit{daemon} \gets \Call{LiveEvaluation}{\mathit{event},
           \mathit{event.asset.predictors},\,\mathit{evaluation}}$
    \State $\mathit{predictors} \gets \Call{AssignMetricPredictors}{\mathit{daemon},\mathit{event}}$
    \BlankLine
    \ForAll{$(\mathit{pName},\mathit{runID}) \in
            \Call{RunPredictor}{\mathit{daemon},\,\mathit{predictors}}$}
        \ForAll{$m \in \mathit{predictors}[\,\mathit{pName}\,].\textit{metrics}$}
            \State $\mathit{data} \gets 
                   \mathit{daemon.predictors}[\,\mathit{pName}\,]
                   .\Call{getMetricData}{\mathit{runID},\,m}$
            \If{$\mathit{data} \neq \varnothing$}
               \State $c_{\mathrm{ok}} \gets c_{\mathrm{ok}} + 1$
               \State $\Call{SetMetricData}{\mathit{daemon},\,m,\,\mathit{pName},\,\mathit{data}}$
            \EndIf
        \EndFor
    \EndFor
    \BlankLine
    \If{$c_{\mathrm{ok}} > 0$}
       \State $\mathit{evaluation.status} \gets \textsc{Finished}$
    \Else
       \State $\mathit{evaluation.status} \gets \textsc{Invalid}$
    \EndIf
    \State $\Call{save}{\mathit{evaluation}} \;\;|\;\;
            \Call{save}{\mathit{daemon}}$
    \State \Return $\mathit{evaluation}$
  \EndFunction
\end{algorithmic}
\end{algorithm}


\begin{algorithm}[h]
\caption{\enskip Assigning predictors to active metrics (method of \textsc{LiveEvaluation}).}\label{alg:assign}
\begin{algorithmic}
  \Function{AssignMetricPredictors}{$\mathit{daemon},\,\mathit{event}$}
    \State $\mathit{Q} \gets \varnothing$ \Comment{Map: predictor $\mapsto$ (runID, metric-list)}
    \ForAll{$m \in \mathit{daemon.active\_metrics}$}
        \ForAll{$(\pi, s) \in 
                 \Call{ScorePredictors}{m,\,\mathit{daemon.predictors},\,\mathit{event}}$}
            \If{$\pi.\textit{active} \land \pi.\textit{name} \notin \operatorname{dom}(\mathit{Q})$}
               \State $\rho \gets \pi.\Call{newPrediction}{\mathit{event}}$
               \State $\mathit{Q}[\pi.\textit{name}] \gets (\rho,\,[\,])$
            \EndIf
            \State $\pi.\Call{activateMetric}{m,\;\mathit{Q}[\pi.\textit{name}].\rho}$
            % \State $\Call{append}{\mathit{Q}[\pi.\textit{name}].\textit{metrics}, \Call{AddMetric}{\mathit{daemon},\,m,\,\pi.\textit{name},\,s}}$
        \EndFor
    \EndFor
    \State \Return $\mathit{Q}$
  \EndFunction
\end{algorithmic}
\end{algorithm}



\FloatBarrier